%% kapa: exam=2 round=34 year=2023 period=1 subject=appraisal_practice subject_ko=감정평가실무 seq=4 number=4 points=10
%% source: Q-Net 감정평가사 기출문제 (cst00309) artlSeq=5212168 file=raw/1008/5212168/1교시 23년도 감정평가사 2차시험_감정평가실무.pdf
%% conversion: kordoc 4.24.51 -> tools/gen2.py (auto); review=proofread-gpt-sol; flags=none
\begin{longquestion}{【문제 4】}
다음과 같이 도로개설사업에 편입되고 남은 토지(잔여지)가 있다. 잔여지 손실보상 기준을 약술하고, 적정한 보상액을 구하시오. (10점) \par
< 자료 1 > 기본적 사항 \par
1. 토지면적: 2,000m² (편입면적: 1,700m²) \par
2. 토지특성: 소로한면 / 세장형 / 평지 \par
3. 편입토지 보상평가액(평균): @600,000원/m² \par
< 자료 2 > 토지특성 자료 \par
1. 도로접면 \par
{\small\centering
\resizebox{\ifdim\width>\linewidth\linewidth\else\width\fi}{!}{%
\begin{tabular}{|c|c|c|c|c|}
\hline
\makecell{구분} & \makecell{소로한면} & \makecell{세로(가)} & \makecell{세로(불)} & \makecell{맹지} \\
\hline
\makecell{소로한면} & \makecell{1.00} & \makecell{0.90} & \makecell{0.80} & \makecell{0.65} \\
\hline
\makecell{세로(가)} & \makecell{1.10} & \makecell{1.00} & \makecell{0.89} & \makecell{0.72} \\
\hline
\makecell{세로(불)} & \makecell{1.25} & \makecell{1.12} & \makecell{1.00} & \makecell{0.81} \\
\hline
\makecell{맹지} & \makecell{1.53} & \makecell{1.38} & \makecell{1.23} & \makecell{1.00} \\
\hline
\end{tabular}}\par}

2. 형상 \par
{\small\centering
\resizebox{\ifdim\width>\linewidth\linewidth\else\width\fi}{!}{%
\begin{tabular}{|c|c|c|c|c|}
\hline
\makecell{구분} & \makecell{가장(정방)형} & \makecell{세장형} & \makecell{사다리형} & \makecell{부정형} \\
\hline
\makecell{가장(정방)형} & \makecell{1.00} & \makecell{0.98} & \makecell{0.97} & \makecell{0.92} \\
\hline
\makecell{세장형} & \makecell{1.02} & \makecell{1.00} & \makecell{0.99} & \makecell{0.94} \\
\hline
\makecell{사다리형} & \makecell{1.03} & \makecell{1.01} & \makecell{1.00} & \makecell{0.95} \\
\hline
\makecell{부정형} & \makecell{1.08} & \makecell{1.06} & \makecell{1.05} & \makecell{1.00} \\
\hline
\end{tabular}}\par}

3. 고저 \par
{\small\centering
\resizebox{\ifdim\width>\linewidth\linewidth\else\width\fi}{!}{%
\begin{tabular}{|c|c|c|c|c|}
\hline
\makecell{구분} & \makecell{평지} & \makecell{완경사} & \makecell{급경사} & \makecell{저지} \\
\hline
\makecell{평지} & \makecell{1.00} & \makecell{0.98} & \makecell{0.92} & \makecell{0.90} \\
\hline
\makecell{완경사} & \makecell{1.02} & \makecell{1.00} & \makecell{0.94} & \makecell{0.92} \\
\hline
\makecell{급경사} & \makecell{1.08} & \makecell{1.06} & \makecell{1.00} & \makecell{0.98} \\
\hline
\makecell{저지} & \makecell{1.11} & \makecell{1.08} & \makecell{1.02} & \makecell{1.00} \\
\hline
\end{tabular}}\par}

< 자료 3 > 기타 자료 \par
1. 잔여지 특성: 맹지 / 부정형 / 저지 \par
2. 잔여지는 도로사업의 시행으로 맹지가 됨은 물론 일반적인 경우와는 달리 기존의 마을과 단절되어 마을과 연결되는 통로의 개설이 필요한 상태로 사업시행자도 이를 인정하고 있음 \par
3. 마을과 잔여지를 연결하는 통로[세로(가)]의 개설 비용: 150,000,000원/식(제반 부대비용을 포함) \par
\end{longquestion}
